% ============================================================================
% Fichier  : partie6/P06_C16_mesurer_articuler.tex
% Titre    : Mesurer la Conformité, Articuler les Normes
% Auteur   : MINKA MI NGUIDJOÏ Thierry Emmanuel
% Version  : 1.0 -- novembre 2026
% Corrigés : corriges/Corr_C16.tex
% ============================================================================
% Sources : auditologie (hal-05701007, hal-05701270) ; Shannon (1948) ;
% Minka Mi Nguidjoï, Massoussi Songo et Bialo (2026, manuscrit en
% préparation, septembre 2026) : champ de la loi (art. 2, 4, 72), tensions
% T1 à T3, matrice en huit étapes, échelle compatibilité / tension /
% incompatibilité, grille de qualification des acteurs, quatre propositions.
% Les éléments du manuscrit sont présentés comme une piste de recherche en
% cours de l'auteur et de ses coauteurs.
% ============================================================================

\chapter{Mesurer la Conformité, Articuler les Normes}
\label{chap:recherche}
\chaptermark{Mesurer et articuler}

\epigraph{%
    « Ce que l'on ignore à cause des faits se cherche.\\
    Ce que l'on ignore à cause des textes se signale. »%
}{Principe du cours}

\niveauChapitre{Recherche}{%
    Tout le cours, en particulier les Chapitres~\ref{chap:audit},
    \ref{chap:loi} et~\ref{chap:programme}
}

\begin{boxobjectifs}
\begin{itemize}
    \item Mesurer l'état de conformité d'une organisation par $\mu$, $U$
          et $H$, et en suivre l'évolution.
    \item Distinguer et mesurer séparément l'incertitude de preuve et
          l'incertitude de critère.
    \item Présenter les pistes de recherche actuelles de l'auteur et de ses
          coauteurs sur l'articulation de la loi n°~2024/017 avec la chaîne
          pénale.
    \item Appliquer une matrice d'analyse des tensions normatives à un cas.
    \item Formuler une question de recherche sur la mesure de la
          conformité ou l'articulation des normes.
\end{itemize}
\end{boxobjectifs}

\bigskip

% ============================================================================
\section*{Avant de commencer : « dix inconnus, est-ce notre faute ? »}
% ============================================================================

Bafoussam, vendredi 25~septembre~2026. Le conseil d'administration de NYANGA
CRÉDIT a lu la synthèse~: vingt-huit écarts, huit exigences satisfaites,
dix inconnues. Un administrateur demande~: « Ces dix inconnues, est-ce notre
faute~? » Mme~BELLO répond que quatre le sont~: des pièces que
l'établissement n'a pas pu produire. Les six autres ne dépendent de
personne dans la salle~: elles attendent un décret, une liste de pays, un
référentiel, ou la décision d'un juge sur le sens d'un article.

La réponse ouvre deux questions de recherche. La première est de mesure~:
peut-on séparer, dans un rapport, ce que l'organisation ignore de ce que
le droit lui-même ne dit pas encore~? La seconde est d'articulation~: que
faire quand deux textes applicables à la même donnée se rencontrent~?

\begin{boxalert}
\textbf{Violation typique.} Imputer à l'organisation auditée toute
l'incertitude d'un audit, y compris celle qui vient de textes absents ou
contradictoires.

\smallskip
\textbf{Règle~:} l'incertitude se décompose selon sa source. Le rapport
distingue ce qui manque dans les preuves de ce qui manque dans les
critères, et il présente les secondes comme une information sur l'état du
droit, pas comme un défaut de l'organisation.
\end{boxalert}

% ============================================================================
\section{Mesurer l'État de Conformité}
\label{sec:c16-mesurer}
% ============================================================================

\begin{boxdef}
Pour un ensemble de $n$ exigences cotées~: $\mu = n_O/n$ (taux de
conformité), $U = n_K/n$ (taux d'incertitude), $H = -\sum_c p_c \log_2 p_c$
(entropie de la répartition des cotes, en bits, entre~0 et
$\log_2 3 \approx 1{,}585$)~\cite{minka2026nek, shannon1948}.
\end{boxdef}

\begin{boxlecture}{NYANGA CRÉDIT, deux instantanés}
\begin{tableaucours}{L{3.4cm}cccccc}{%
    \tableauHeader{Date} &
    \tableauHeader{$n$} &
    \tableauHeader{O / N / K} &
    \tableauHeader{$\mu$} &
    \tableauHeader{$U$} &
    \tableauHeader{$H$} &
    \tableauHeader{}
}
6~juillet~2026 & 8 & 2 / 2 / 4 & 0,25 & 0,50 & 1,50 & premier jour \\
25~septembre~2026 & 46 & 8 / 28 / 10 & 0,17 & 0,22 & 1,35 & rapport \\
\end{tableaucours}

\medskip
\textbf{Lecture.} Entre les deux dates, $U$ baisse de 0,50 à 0,22~:
l'audit a appris. $\mu$ baisse aussi, de 0,25 à 0,17, sans que
l'établissement ait régressé~: les \NEK{} du premier jour se sont résolus
surtout en \NON{}. Une baisse de $\mu$ au cours d'un audit est donc souvent
un signe de connaissance, pas de dégradation. Les deux instantanés ne
portent pas sur le même ensemble d'exigences~; la comparaison n'est
rigoureuse que sur les huit exigences communes.
\end{boxlecture}

La dernière remarque est une première question de recherche~: comment
comparer deux états de conformité mesurés sur des ensembles d'exigences
différents, ou découpés différemment~? Les mesures dépendent de la
granularité~; un domaine découpé en dix exigences pèse plus qu'un domaine
découpé en quatre.

% ============================================================================
\section{Deux Sortes d'Incertitude}
\label{sec:c16-deux-nek}
% ============================================================================

Le cours a distingué tout au long deux sources de \NEK{}~:
\begin{itemize}
    \item le \textbf{\NEK{} de preuve}~: le critère est clair, la pièce
          manque~; il se lève en cherchant (audit du prestataire, examen des
          sauvegardes)~;
    \item le \textbf{\NEK{} de critère}~: la pièce existerait, mais la règle
          qui permettrait de la juger n'est pas fixée~; il se lève avec un
          texte (décret, référentiel, liste de pays) ou une décision de
          justice.
\end{itemize}

Le cours propose, à titre de piste, de décomposer le taux d'incertitude en
deux termes~:
\[
U = U_p + U_c, \qquad U_p = \frac{n_{K,p}}{n}, \qquad U_c = \frac{n_{K,c}}{n},
\]
où $n_{K,p}$ et $n_{K,c}$ comptent les \NEK{} de preuve et de critère. À
NYANGA, sur 46 exigences, quatre \NEK{} sont de preuve et six de critère~:
$U_p \approx 0{,}09$ et $U_c \approx 0{,}13$.

$U_p$ mesure ce que l'organisation peut réduire. $U_c$ mesure ce que l'état
du droit laisse ouvert~; mesuré sur plusieurs organisations et plusieurs
années, il pourrait servir d'indicateur de la maturité du cadre normatif
lui-même, et décroître à mesure que les textes d'application paraissent.
Cette décomposition n'est pas encore validée~: elle suppose une règle sûre
pour classer chaque \NEK{}, et certains relèvent des deux sources à la
fois.

% ============================================================================
\section{Une Piste de Recherche Actuelle : la Loi et la Chaîne Pénale}
\label{sec:c16-piste}
% ============================================================================

\begin{boxinfo}
\textbf{Statut de cette section.} Elle présente un travail en cours de
l'auteur de ce cours et de ses coauteurs, M.~Robert Thierry MASSOUSSI SONGO
et M.~Dieudonné BIALO, intitulé « Protection des données personnelles et
chaîne pénale au Cameroun~: analyse des zones de tension normative autour
de la Loi n°2024/017 »~\cite{minka2026chainepenale}. Ce manuscrit n'est pas
encore publié~; ses analyses et ses propositions sont celles des auteurs,
et elles peuvent évoluer. Les cas qu'il étudie sont construits à des fins
d'analyse.
\end{boxinfo}

\subsection{La question}

Les auteurs partent d'un constat~: la loi n°~2024/017 institue un régime
général, mais ses articles~4 et~72 renvoient la sécurité, la défense et la
justice à des textes particuliers ou spécifiques. Leur problématique est
de savoir dans quelle mesure la loi constitue le régime de référence pour
les traitements de la chaîne pénale (police judiciaire, expertise,
magistrature), et quelles zones de tension subsistent avec les autres
instruments applicables~: le Code pénal, le Code de procédure pénale, la
loi n°~2010/012, la Convention de Budapest à laquelle le Cameroun est
Partie, et la Convention de Malabo, signée mais non ratifiée.

\subsection{Trois niveaux d'articulation}

Les auteurs distinguent trois situations entre deux normes applicables à
la même opération~:
\begin{itemize}
    \item la \textbf{compatibilité}~: elles s'appliquent simultanément sans
          difficulté~;
    \item la \textbf{tension}~: leur application simultanée crée une
          difficulté qui demande une interprétation ou des garanties~;
    \item l'\textbf{incompatibilité}~: les deux obligations ne peuvent pas
          juridiquement être satisfaites ensemble.
\end{itemize}

Cette distinction, selon eux, évite de transformer toute difficulté
d'articulation en conflit de normes.

\subsection{La matrice d'analyse}

\begin{tableaucours}{cL{3.2cm}L{8.0cm}}{%
    \tableauHeader{N°} &
    \tableauHeader{Étape} &
    \tableauHeader{Question}
}
1 & Applicabilité & Chaque norme est-elle applicable~? \\
2 & Qualification & Quelle est la qualification juridique de l'opération~? \\
3 & Concurrence & Les normes portent-elles sur la même opération~? \\
4 & Conciliabilité & Imposent-elles réellement des obligations
simultanées~? \\
5 & Interprétation & Peut-on les concilier par interprétation~? \\
6 & Qualification de la tension & Compatibilité, tension ou
incompatibilité~? \\
7 & Mécanisme & Quel mécanisme de résolution est envisageable~? \\
8 & Proposition & Quelle solution proposer~? \\
\end{tableaucours}
\vspace{-0.8em}
\begin{center}{\small\itshape Tableau~16.1. La matrice des auteurs (d'après~\cite{minka2026chainepenale})}\end{center}

\subsection{Trois zones de tension}

\begin{tableaucours}{cL{3.0cm}L{8.2cm}}{%
    \tableauHeader{Zone} &
    \tableauHeader{Nature} &
    \tableauHeader{Description}
}
T1 & Temporelle & Les données collectées avant la loi, puis conservées,
analysées ou produites en justice après son entrée en vigueur~: quelle loi
régit chaque opération~? \\
T2 & Procédurale & L'information « sans délai » de la personne
(article~22) et le secret de l'enquête et de l'information judiciaire
(articles~102 et~154 du Code de procédure pénale) \\
T3 & Transfrontalière & La conservation rapide de la Convention de Budapest
(article~29) et l'autorisation préalable des articles~19 et~32, sous
réserve de l'article~72 \\
\end{tableaucours}
\vspace{-0.8em}
\begin{center}{\small\itshape Tableau~16.2. Les trois zones de tension (d'après~\cite{minka2026chainepenale})}\end{center}

Pour la zone temporelle, les auteurs privilégient une solution où la
légalité de la collecte est régie par la loi en vigueur au moment de la
collecte, en distinguant ensuite conservation et traitement ultérieur.
Pour la zone transfrontalière, ils recommandent de distinguer
conservation, préservation, production, transmission et transfert, qui
peuvent relever de régimes différents. Leur conclusion est que ces zones
ne constituent pas, à ce stade, des incompatibilités irréductibles~: leur
résolution appelle, selon les cas, une clarification jurisprudentielle,
réglementaire ou législative.

\subsection{Qualifier les acteurs opération par opération}

\begin{tableaucours}{L{2.4cm}L{9.8cm}}{%
    \tableauHeader{Acteur} &
    \tableauHeader{Qualification possible selon les auteurs}
}
Magistrat & Responsable du traitement lorsqu'il détermine les finalités et
les principaux éléments du traitement, sous réserve des circonstances \\
OPJ & Responsable du traitement ou autre acteur habilité, selon le cadre
institutionnel, les finalités et son autonomie de décision \\
Expert & Sous-traitant possible lorsqu'il traite pour le compte d'un
responsable sans déterminer lui-même les finalités \\
\end{tableaucours}
\vspace{-0.8em}
\begin{center}{\small\itshape Tableau~16.3. La grille fonctionnelle (d'après~\cite{minka2026chainepenale})}\end{center}

\subsection{Quatre propositions}

Les auteurs formulent quatre propositions~: clarifier le champ
d'application de la loi aux traitements judiciaires~; distinguer les
opérations dans les échanges transfrontaliers~; prévoir un mécanisme qui
permette au magistrat compétent d'apprécier le moment et les modalités de
l'information quand elle risque de compromettre une enquête~; élaborer une
doctrine sur la temporalité des traitements. Ils annoncent des études
complémentaires sur l'effet d'une éventuelle ratification de la Convention
de Malabo, sur la jurisprudence à venir, et sur les textes d'application de
la loi lorsqu'ils seront publiés.

% ============================================================================
\section{Deux Échelles à Trois Niveaux}
\label{sec:c16-echelles}
% ============================================================================

Le cours rapproche, à titre de question ouverte, deux instruments de
l'auteur. Les cotes \OUI{}, \NON{}, \NEK{} portent sur des \textit{faits}~:
une exigence est-elle satisfaite, au vu des preuves~? L'échelle
compatibilité, tension, incompatibilité porte sur des \textit{normes}~:
deux textes peuvent-ils être satisfaits ensemble~?

Une hypothèse de recherche se dessine~: une exigence qui relève d'une zone
de tension non résolue produit, dans un audit, un \NEK{} de critère. La
compatibilité permettrait de coter normalement~; l'incompatibilité
obligerait à choisir un texte et à le dire~; la tension laisserait un
\NEK{} tant qu'aucune interprétation ne s'impose. Si cette hypothèse se
vérifiait, le taux $U_c$ d'un audit pourrait se rattacher aux zones de
tension recensées par l'analyse juridique. Rien de cela n'est encore
établi~; c'est un programme de travail.

% ============================================================================
\section{Questions Ouvertes}
\label{sec:c16-questions}
% ============================================================================

\begin{itemize}
    \item \textbf{La granularité.} Comment normaliser le découpage des
          exigences pour comparer des audits entre eux~?
    \item \textbf{La pondération.} Toutes les exigences ne pèsent pas
          autant pour les personnes. Faut-il pondérer $\mu$ et $U$ par la
          gravité, et comment éviter que la pondération ne devienne
          arbitraire~?
    \item \textbf{La classification des \NEK{}.} Quelle règle permet de
          classer sûrement un \NEK{} en preuve ou en critère~?
    \item \textbf{La maturité normative.} Le taux $U_c$ mesuré sur plusieurs
          organisations permet-il de suivre la mise en place des textes
          d'application~?
    \item \textbf{L'articulation des normes.} La matrice des auteurs
          peut-elle être outillée pour guider l'auditeur devant une exigence
          prise entre deux textes~?
\end{itemize}

\begin{boxinfo}
\textbf{Construire un sujet de mémoire.}
\begin{enumerate}
    \item Choisir un objet~: une organisation, un secteur, un type de
          traitement, une zone de tension.
    \item Fixer l'ensemble des exigences et la règle de cotation.
    \item Recueillir des cas réels anonymisés, avec l'accord des
          organisations, ou des cas construits et documentés.
    \item Mesurer $\mu$, $U_p$, $U_c$ et $H$, ou appliquer la matrice des
          tensions.
    \item Confronter les résultats à un jugement d'experts.
    \item Dire les limites~: taille de l'échantillon, subjectivité des
          cotes, état des textes à la date de l'étude.
\end{enumerate}
\end{boxinfo}

% ============================================================================
\section{Ce que la Mesure demande à chaque Métier}
\label{sec:c16-metiers}
% ============================================================================

\begin{boxmetier}{Auditeur de la protection des données}
Classer chaque \NEK{} en preuve ou en critère, publier $U_p$ et $U_c$ à
côté de $\mu$, et dire dans le rapport ce qui ne dépend pas de
l'organisation.
\end{boxmetier}

\begin{boxmetier}{Juriste conformité}
Tenir la liste des zones de tension qui touchent l'organisation, avec,
pour chacune, l'interprétation retenue et sa justification, en attendant
les textes ou la jurisprudence.
\end{boxmetier}

\begin{boxmetier}{Responsable de traitement}
Suivre d'un audit à l'autre l'évolution de $U_p$, qui mesure ce que
l'organisation peut encore apprendre sur elle-même.
\end{boxmetier}

\begin{boxresume}
\begin{itemize}
    \item $\mu$, $U$ et $H$ décrivent l'état de conformité~; une baisse de
          $\mu$ au cours d'un audit traduit souvent une connaissance
          nouvelle.
    \item Le cours propose de décomposer $U$ en incertitude de preuve et
          incertitude de critère~; à NYANGA, $U_p \approx 0{,}09$ et
          $U_c \approx 0{,}13$.
    \item L'auteur et ses coauteurs étudient, dans un travail en cours, les
          tensions temporelle, procédurale et transfrontalière entre la loi
          et la chaîne pénale, avec une matrice en huit étapes.
    \item Ils distinguent compatibilité, tension et incompatibilité, et
          qualifient les acteurs opération par opération.
    \item Le rapprochement de cette échelle et des cotes \OUI{}, \NON{},
          \NEK{} est une question de recherche ouverte.
\end{itemize}
\end{boxresume}

\begin{boxinfo}
\textbf{Points de contrôle.} Vous avez compris ce chapitre si vous savez~:
\begin{itemize}
    \item calculer $U_p$ et $U_c$ pour un ensemble de cotes~;
    \item appliquer la matrice des tensions à un cas simple~;
    \item formuler une question de recherche et ses limites.
\end{itemize}
\end{boxinfo}

\begin{boxexo}[title={\ding{45}\quad Exercice~16.1. Calculer \niveaubase}]
Pour un audit de 20 exigences~: 6 \OUI{}, 9 \NON{}, 3 \NEK{} de preuve et
2 \NEK{} de critère. Calculez $\mu$, $U$, $U_p$, $U_c$ et $H$.
\end{boxexo}

\begin{boxexo}[title={\ding{45}\quad Exercice~16.2. Classer les NEK \niveaumoyen}]
Classez en preuve ou en critère les \NEK{} suivants, rencontrés dans le
cours~: (a) l'exécution des mesures par l'éditeur du logiciel bancaire~;
(b) le niveau de protection de la Côte d'Ivoire~; (c) la licéité du
profilage au regard de l'article~65~; (d) le pays de traitement des
selfies~; (e) l'application de l'article~22 en mars~2025.
\end{boxexo}

\begin{boxexo}[title={\ding{45}\quad Exercice~16.3. La matrice \niveaumoyen}]
Appliquez les huit étapes de la matrice à la tension T2 dans l'affaire
NYANGA~: la demande de l'OPJ du 20~mars~2025 et l'article~22. À quelle
étape la situation change-t-elle en février~2026~?
\end{boxexo}

\begin{boxexo}[title={\ding{45}\quad Exercice~16.4. Un sujet de mémoire \niveauexpert}]
Choisissez l'une des questions ouvertes de la section~\ref{sec:c16-questions}
et rédigez une proposition de mémoire d'une page~: question, hypothèse,
données, méthode, critère de validation, limites.
\end{boxexo}

\bigskip

\begin{boxinfo}
\textbf{Pour aller plus loin.}
\begin{itemize}
    \item T.~E.~Minka Mi Nguidjoï, « Vers l'auditologie »~\cite{minka2026auditologie}.
    \item T.~E.~Minka Mi Nguidjoï, sur l'instrument trichotomique~\cite{minka2026nek}.
    \item T.~E.~Minka Mi Nguidjoï, R.~T.~Massoussi Songo et D.~Bialo,
          manuscrit en préparation~\cite{minka2026chainepenale}.
\end{itemize}
\end{boxinfo}

\begin{boxinfo}
\textbf{Fin du cours.} L'affaire a commencé par une lettre, le lendemain
d'une échéance~: « Sommes-nous en règle~? » Elle s'achève sur un rapport
qui répond non sur vingt-huit points, oui sur huit, et qui dit pourquoi il
ne peut pas répondre sur dix. Entre les deux, l'audit a recensé, apprécié,
vérifié, coté et recommandé. La conformité n'est pas un état que l'on
déclare~; c'est une preuve que l'on construit, et dont on dit les limites.
\end{boxinfo}
